\documentclass[12pt,a4paper]{article}

\usepackage{fontspec}
\usepackage{amsmath, amssymb, amstext}
\usepackage{graphicx}
\usepackage{geometry}
\usepackage{xcolor}
\usepackage{fancyhdr}

\geometry{margin=1in}

\pagestyle{fancy}
\setlength{\headheight}{15pt}

\lhead{کوییز درس فیزیک مدرن}
\chead{دانشکده فیزیک دانشگاه صنعتی خواجه نصیرالدین طوسی}
\rhead{۱۴۰۵/۰۳/۲۶}

\usepackage{xepersian}
\settextfont{B Nazanin}

\begin{document}
	
	\begin{tabular}{|p{7cm}|p{7cm}|}
		\hline
		\textbf{نام و نام خانوادگی:} & \textbf{شماره دانشجویی:} \\
		\hline
	\end{tabular}
	
	\vspace{10pt}
	
	\section*{سوالات}
	
	\begin{enumerate}
		
		\item 
		تکانه ذرات زیر را محاسبه کنید:
		
		\begin{enumerate}
			\item[(الف)] یک فوتون گاما با انرژی
			\[
			8.0\,\mathrm{MeV}
			\]
			
			\item[(ب)] یک فوتون پرتو ایکس با انرژی
			\[
			18\,\mathrm{keV}
			\]
			
			\item[(پ)] یک فوتون فروسرخ با طول موج
			\[
			1.5\,\mu\mathrm{m}
			\]
			
			\item[(ت)] یک فوتون امواج رادیویی با فرکانس
			\[
			120\,\mathrm{MHz}
			\]
		\end{enumerate}
		
		\item
		نوری با طول موج \(280.0\,\mathrm{nm}\) به سطح یک فلز تابیده می شود. طول موج آستانه این فلز برابر \(360.0\,\mathrm{nm}\) است.
		
		با استفاده از معادله اثر فوتوالکتریک اینشتین، پتانسيل توقف \lr{(Stopping Potential)} را محاسبه کنید.
		
	\end{enumerate}
	
	\vspace{10pt}
	\textbf{موفق باشید.}
	
\end{document}